% Generated draft: requires mathematical and source review.
\documentclass[11pt]{article}
\usepackage[margin=25mm]{geometry}
\usepackage{fontspec}
\setmainfont{Latin Modern Roman}
\usepackage{amsmath,amssymb,mathtools,unicode-math}
\setmathfont{Latin Modern Math}
\usepackage{enumitem,needspace,xurl,hyperref}
\setlength{\parindent}{0pt}
\setlength{\parskip}{5pt}
\setlength{\emergencystretch}{3em}
\newcommand{\N}{\mathbb{N}}
\newcommand{\Z}{\mathbb{Z}}
\newcommand{\Q}{\mathbb{Q}}
\newcommand{\R}{\mathbb{R}}
\newcommand{\C}{\mathbb{C}}
\newcommand{\F}{\mathbb{F}}
\newcommand{\A}{\mathbb{A}}
\newcommand{\PP}{\mathbb{P}}
\newcommand{\E}{\mathbb{E}}
\newcommand{\eps}{\varepsilon}
\newcommand{\dd}{\,\mathrm{d}}
\begin{document}
\section*{UnsolvedMath ID: 30001840}
\textbf{Source problem code:} catalogue entry
\subsection{Definitions and mathematical statement}
Put
\[
\mathcal O=\mathbb Z\!\left[\frac{1+\sqrt5}{2}\right],\qquad f_t(x)=x^5-5x^3+5x+t.
\]
The supplied source considers the one-parameter family \(C_t:y^2=f_t(x)\) and attributes real multiplication by \(\mathcal O\) to its Jacobians. A genus-two curve here must mean the smooth projective model of this affine equation at a parameter for which the polynomial is squarefree. Singular parameters cannot be included as smooth genus-two fibers.

The source does not specify a base field or the intended range of parameters. To define the arithmetic objects conditionally, let \(K\) be a field of characteristic zero and let \(t\in K\) satisfy the squarefreeness condition. Write \(J_t=\operatorname{Jac}(C_t)\), choose an algebraic closure \(\overline K\), and put \(G_K=\operatorname{Gal}(\overline K/K)\). These are auxiliary choices for explaining the question, not a claim that the source quantifies over all such fields and parameters. Real multiplication means an embedding
\[
\iota:\mathcal O\hookrightarrow\operatorname{End}_{\overline K}(J_t).
\]
Whether the intended embedding is defined over \(K\), and what compatibility with the polarization is required, are not specified in the supplied record.

For a positive integer \(N\), define
\[
J_t[N](\overline K)=\{P\in J_t(\overline K):NP=0\},\qquad
\rho_{t,N}:G_K\longrightarrow\operatorname{Aut}_{\mathbb Z/N\mathbb Z}\bigl(J_t[N](\overline K)\bigr)
\]
by the natural Galois action. For a prime \(\ell\), the \(\ell\)-adic Tate module and its action are
\[
T_\ell(J_t)=\varprojlim_n J_t[\ell^n](\overline K),\qquad
\rho_{t,\ell^\infty}:G_K\longrightarrow\operatorname{Aut}_{\mathbb Z_\ell}(T_\ell(J_t)),
\]
where the transition maps are multiplication by \(\ell\). An alternative possible meaning of the source's full Tate-module question is the action on \(\prod_{\ell}T_\ell(J_t)\); the source does not distinguish these targets.

\textbf{Source question.} Inside the two-dimensional real-multiplication moduli space referred to by the source, determine in what sense the image of this one-dimensional family is special, and determine which level structure the family fixes. Determine whether an exceptional positive integer \(N\) can force the mod-\(N\) Galois image to be smaller than expected, and determine the image on the intended Tate module.

A fully specified arithmetic question still requires a choice of base field, generic fiber or specialization, parameter quantifiers, and comparison group for the expected image. The comparison group must account for the intended real-multiplication and polarization data and their fields of definition. The source supplies none of these choices explicitly. In particular, the displayed conditional representations do not resolve whether the question concerns arithmetic Galois images, geometric monodromy, or both; no particular interpretation or expected maximal group is asserted here.
\subsection{Short English statement}
For the smooth genus-two members of \(y^2=x^5-5x^3+5x+t\), whose Jacobians are stated to have real multiplication by \(\mathbb Z[(1+\sqrt5)/2]\), describe the family's special position in the two-dimensional real-multiplication moduli space, its fixed level structure, and its torsion and Tate-module Galois images. The base field, parameter quantifiers, expected image, and meaning of the Tate-module target remain unspecified.
\subsection{Sources}
\begin{itemize}
\item Local UnsolvedMath registry, supplied complete record: permanent ID \(30001840\), source problem code \texttt{catalogue entry}, title \emph{Galois Images in Genus-Two Real-Multiplication Families}, in the Oberwolfach Reports Open Problems collection. No individual proposer is supplied; the proposed year is recorded as 2011.
\item Supplied authoritative report citation and literature-source label ``original OWR report'': \emph{Explicit Methods in Number Theory (2012)}, Oberwolfach Reports, pp.~1971--2044, DOI \texttt{10.4171/owr/2011/35}. \url{https://doi.org/10.4171/owr/2011/35}. The collection describes these records as drawn from mathematical workshop reports and their open-problem discussions.
\item Supplied literature-source label ``integral and adelic open-image results'': \url{https://arxiv.org/abs/1508.06426}. No authors, title, or precise applicable theorem are supplied for this reference.
\item Literature assessment checked at \texttt{2026-08-21}: the supplied status is \texttt{partially\_solved}. The supplied assessment says that Galois images are determined in several low-dimensional and sufficiently generic real-multiplication families, while exceptional-level and full-Tate-module questions remain open in the stated generality. This is supplied dated curation derived from the OWR collection, not an independently verified literature assessment; its time-sensitive status and cited claims require independent checking.
\item The registry labels the statement \texttt{corrected\_verified} and reports that its mathematical content was checked against the official report, with notation, context, and question form normalized without changing the source-level claim. This records the supplied verification assertion; no independent source visit or mathematical verification is claimed here.
\end{itemize}
\end{document}
